% Marian antiphons, sung after Mass; \input by every booklet.
% Liturgical year from Advent; each antiphon in the simple, then the solemn tone.
% ---------------------------------------------------------------------------
\antiphonsheading

% ---------------------------------------------------------------------------
\chant{Alma Redemptoris Mater}{tonus V}
{\centering\small\rubr{Od I Niedzieli Adwentu do święta Ofiarowania Pańskiego (2 lutego).}\par}
\par\smallskip{\centering\small\itshape ton prosty\par}
\lilypondfile{../shared/scores/alma-simplex.ly}
\par\smallskip{\centering\small\itshape ton uroczysty\par}
\lilypondfile{../shared/scores/alma.ly}
\begin{tekst}
\w{Alma Redemptóris Mater, quae pérvia caeli porta manes,}
  {Święta Matko Odkupiciela, Bramo niebios zawsze otwarta,}
\w{et stella maris, succúrre cadénti, súrgere qui curat pópulo:}
  {i Gwiazdo morza, wspomóż upadający lud, który pragnie powstać.}
\w{tu quae genuísti, natúra miránte, tuum sanctum Genitórem,}
  {Ty, która ku zdumieniu natury zrodziłaś swego świętego Stwórcę,}
\w{Virgo prius ac postérius, Gabriélis ab ore sumens illud Ave,
   peccatórum miserére.}
  {Dziewico przedtem i potem, z ust Gabriela przyjmująca owo „Ave”,
   zmiłuj się nad grzesznikami.}
\end{tekst}

\clearpage
% ---------------------------------------------------------------------------
\chant{Ave Regina caelorum}{tonus VI}
{\centering\small\rubr{Od Ofiarowania Pańskiego (2 lutego) do Wielkiej Środy.}\par}
\par\smallskip{\centering\small\itshape ton prosty\par}
\lilypondfile{../shared/scores/ave-regina-simplex.ly}
\par\smallskip{\centering\small\itshape ton uroczysty\par}
\lilypondfile{../shared/scores/ave-regina.ly}
\begin{tekst}
\w{Ave, Regína caelórum, ave, Dómina Angelórum:}
  {Witaj, Królowo niebios, witaj, Pani aniołów,}
\w{salve, radix, salve, porta, ex qua mundo lux est orta:}
  {witaj, Korzeniu, witaj, Bramo, z której Światłość wzeszła dla świata.}
\w{Gaude, Virgo gloriósa, super omnes speciósa:}
  {Raduj się, Panno chwalebna, piękniejsza nad wszystkie;}
\w{vale, o valde decóra, et pro nobis Christum exóra.}
  {bądź pozdrowiona, o najpiękniejsza, i przyczyń się za nami do Chrystusa.}
\end{tekst}

\clearpage
% ---------------------------------------------------------------------------
\chant{Regina caeli}{tonus VI}
{\centering\small\rubr{W okresie wielkanocnym, od Niedzieli Zmartwychwstania
do Zesłania Ducha Świętego.}\par}
\par\smallskip{\centering\small\itshape ton prosty\par}
\lilypondfile{../shared/scores/regina-caeli-simplex.ly}
\par\smallskip{\centering\small\itshape ton uroczysty\par}
\lilypondfile{../shared/scores/regina-caeli.ly}
\begin{tekst}
\w{Regína caeli, laetáre, allelúja:}
  {Królowo nieba, wesel się, alleluja,}
\w{quia quem meruísti portáre, allelúja:}
  {bo Ten, któregoś nosić zasłużyła, alleluja,}
\w{resurréxit, sicut dixit, allelúja:}
  {zmartwychwstał, jak powiedział, alleluja.}
\w{ora pro nobis Deum, allelúja.}
  {Módl się za nami do Boga, alleluja.}
\end{tekst}

\clearpage
% ---------------------------------------------------------------------------
\chant{Salve Regina}{tonus V (ton prosty), tonus I (ton uroczysty)}
{\centering\small\rubr{Od uroczystości Najświętszej Trójcy (po Zesłaniu Ducha Świętego)
do Adwentu.}\par}
\par\smallskip{\centering\small\itshape ton prosty\par}
\lilypondfile{../shared/scores/salve-regina-simplex.ly}
\par\smallskip{\centering\small\itshape ton uroczysty\par}
\lilypondfile{../shared/scores/salve-regina.ly}
\begin{tekst}
\w{Salve, Regína, mater misericórdiae;}
  {Witaj, Królowo, Matko miłosierdzia,}
\w{vita, dulcédo et spes nostra, salve.}
  {życie, słodyczy i nadziejo nasza, witaj!}
\w{Ad te clamámus, éxsules fílii Hevae.}
  {Do Ciebie wołamy wygnańcy, synowie Ewy;}
\w{Ad te suspirámus, geméntes et flentes in hac lacrimárum valle.}
  {do Ciebie wzdychamy, jęcząc i płacząc na tym łez padole.}
\w{Eia ergo, advocáta nostra, illos tuos misericórdes óculos
   ad nos convérte.}
  {Przeto, Orędowniczko nasza, one miłosierne oczy Twoje
   na nas zwróć,}
\w{Et Jesum, benedíctum fructum ventris tui,
   nobis post hoc exsílium osténde.}
  {a Jezusa, błogosławiony owoc żywota Twojego,
   po tym wygnaniu nam okaż.}
\w{O clemens, o pia, o dulcis Virgo María.}
  {O łaskawa, o litościwa, o słodka Panno Maryjo!}
\end{tekst}
